Ontstaan op veengronden,
waar moerassen stonden

tot men ze drooglegde
en wegen aanlegde,

verrezen bouwwerken
met fiere kenmerken.

‘Het land waar ik door waad’,
in de volksmond ‘Landkwaad’,

noemde hij het landgoed,
waar hij werd opgevoed

door een strenge regent
die zijn va had gekend.

Zijn ouders gingen dood
tijdens de hongersnood

door een belegering
van hun oude vesting.

Deze onzekerheid
zorgde voor zijn beleid

om zijn ‘Landkwaad’
te vrijwaren van kwaad.

Chaos werd zo structuur.
Zoals de buitenmuur

die hij liet aanleggen
na zijn rijk te dreggen.

Deze moest afschrikken,
maar deed hem verstikken

in zijn isolement.
Hij had angst niet herkend

in zijn vele soorten.
Hij sloot vaak zijn poorten

als hij zich zwak voelde.
Wat meest op niets stoelde.

Geen gevaar van buiten,
maar het zich afsluiten

van vrees dieper gelegen,
hield zijn liefde tegen.

De muur bood bescherming.
Elke herinnering

aan zijn vele trauma’s
werd vervaagd door de waas

van uiterlijke kracht.
Verdronken in de gracht,

die zijn Donjon omsloot
waar hij zielsrust genoot.

Waar zijn pure wezen
zijn zachtheid kon lezen.

Dit centrum was de plek
waar geen vrees of gebrek

onophoudelijk streed
tegen zijn hartekreet.

Een fikse woontoren,
gevuld met folklore,

naast de havezate
waar zijn ouders zaten.

Daarvoor de binnenplaats.
Een stal en een werkplaats

omringden het voorhof.
Wat men er niet aantrof

waren opslagschuren.
Uit de bange uren

van zijn ouders geleerd,
lagen gehergroepeerd

alle noodvoorraden
klaar om uit te laden.
De gracht hield water vast.
Emotionele last

was er in opgehoopt.
In pijn en angst gedoopt.

Aan de rand van ‘Landkwaad’
werd een kroeg uitgebaat.

Ook als herberg gebruikt
en als bordeel misbruikt.

Rondom de buitenmuur
lag de natte natuur

drooggelegd voor landbouw
en traditiegetrouw

werden alle wouden
als revier gehouden

en als barrière
voor elke affaire

uit de buitenwereld.
In zijn rijk gezeteld

regelde een wachtpoort
door wie hij werd gestoord.

Alle toegangspaden,
die gasten betraden

waren smal, kronkelig,
drabbig en schemerig

onder een bladerdak,
Net als de maniak

zijn neurale banen,
leidend tot zijn wanen.

De Franse soldaten
struikelden gelaten

over etiquette.
De anders zo nette

beroepsofficieren
bleken dronken klieren.

Zou de wapenstilstand,
beklonken op het land

van de landheer, voortgaan?
Hij ergerde zich aan

zijn brutale gasten.
die de meid betasten,

het meubilair braken
en zo brutaal spraken.

Dat kon hij verdragen.
Doch het alsmaar plagen

over zijn rode haar,
daarmee was hij snel klaar.

Zijn eigen pretentie
gaf de motivatie

op basis van zelfbeeld.
Werd zijn rol uitgespeeld

of kon hij iets leren
om zijn ziel te weren?

“Stelt zo de Franse kroon,
zijn tucht en eer ten toon?

Verwend met mijn wijnen,
maar feestend als zwijnen?

Of kunnen de Fransen
slechts vechten als ganzen?”

Een onderofficier,
zwaaiend met een glas bier,

wankelt van dronkenschap
ontving het als een grap.

“No, Monsieur le Baron,
dit hele bataljon

die u net heeft onteerd,
heeft de Chausson geleerd

in de stad Marseille!”
De alcoholvrije

onbewogen landheer,
knipperde niet een keer

toen de Franse soldaat
met serieus gelaat

fraai trapte in de lucht.
Men lachte om de klucht.

Behalve de landheer
met zijn gekwetste eer.

“Ik heb koken geleerd,
dus ganzen gefileerd!”.

Ieder hield zijn wafel
toen een mes van tafel

in de landheer zijn hand
ongemerkt was beland

en de soldaat neerstak.
Zo staakte hun bivak.

“Rijkdom is veiligheid”
was het landheers beleid.

Zijn hofmeester bewaakte
of zijn geld opraakte.

De wortel was zijn angst.
Wantrouwen de ontvangst

die de vertakkingen
als voeding ontvingen.

Dysfunctie, controle
en perfectionisme

de rotte bladeren.
Stress in de aderen.

Zonder intimiteit
raakte hij mensen kwijt.

Zijn scheppende onschuld
was al als kind omhuld

middels een huid zo bleek
dat het doorschijnend leek.

Doch Nietzsche zag die geest
als gedaante geweest.

Het leven in schokken
eruit weggetrokken.

Zwetende poriën
door de financiën.

De mond die openviel
en door een klein ventiel

scheetsgewijs ademde.
En angst inboezemde.

De verschrikte ogen
van passie vervlogen
kon hij niet intern zien.
Het ego had voorzien

dat de vestingwallen
niet om konden vallen.

Vele handelaren
verkochten hun waren

op het landgoed zijn markt.
Het werd vaak een infarct

van boeren, landlopers,
marktkramers en kopers.

Als producten centrum
bood het een maximum

aan mogelijkheden.
Vaak werd overschreden

het te krappe budget.
Al werd goed opgelet

door de penningmeester,
mocht plots de hofmeester

niet zelf inkopen doen.
De landheer kon zijn poen

niemand toevertrouwen.
Dus ging uit wantrouwen

de landheer dit keer mee.
Hij bleek snel niet tevree

over het afdingen.
Hij kon meer afdwingen

als hij het zelf voordeed.
Bij de volgende keet

zag hij prachtig breiwerk.
Dus maakte hij zich sterk

en zei: “Wat kost dit vod?”
Hij deed echter zijn bod

aan een jonge boerin
gebogen onderin

de karige toonbank
bestaand uit slechts één plank,

dus ondersteboven.
En ondersteboven

was meteen de landheer
toen de vrouw haar verweer

al opkrabbelend gaf.
Haar vraagprijs was vrij straf.

Moeder was zo’n schoonheid
dat plots zonder respijt

het te hoge bedrag
op de uitstalplank lag.

De hofmeester gaf door;
“Scherp onderhandeld hoor!”

Het contact was gelegd!
Wat wilde Moeder echt,

als dochter van Mara,
de duivelse deva?
